El propósito de este Trabajo de Fin de Grado es acercar la consulta de un
catálogo técnico de vehículos a usuarios sin conocimientos de automoción, de
modo que puedan expresar sus necesidades en lenguaje natural y recibir
recomendaciones fundamentadas en datos reales y verificables, y no en
afirmaciones inventadas por el modelo de lenguaje. Con ese fin se ha diseñado,
implementado y validado \emph{autoAI}, una aplicación web que reúne un asistente
conversacional y un catálogo filtrable clásico sobre el mismo conjunto de datos
del mercado español.

El desarrollo siguió una metodología incremental. El sistema se organiza en
cuatro servicios desacoplados y contenerizados: un proceso de extracción
(\emph{web scraping}) con Scrapy, que recorre la jerarquía del agregador
\texttt{km77.com}; un \emph{backend} en FastAPI que normaliza e ingiere los
datos, expone una API REST y orquesta el asistente; una base de datos relacional
PostgreSQL; y un \emph{frontend} de página única en React y TypeScript. El núcleo
del trabajo es un \emph{pipeline} de inteligencia artificial basado en el patrón
de uso de herramientas (\emph{tool calling}): el modelo de lenguaje interpreta la
petición y redacta la respuesta, pero toda cifra que ofrece procede
obligatoriamente de una consulta a la base de datos, lo que impide por
construcción las alucinaciones.

El sistema resultante es funcional y desplegable. El asistente traduce peticiones
vagas en filtros concretos y devuelve respuestas estructuradas (texto, tablas,
gráficas e imágenes) tomadas de datos reales, junto a un catálogo filtrable y
paginado. Se resolvió la normalización de datos heterogéneos y con información
implícita, en particular la inferencia del nivel de electrificación a partir del
nombre comercial, y se diseñó una ingesta idempotente y tolerante a fallos. Las
funcionalidades críticas se respaldaron con pruebas automatizadas, un flujo de
integración continua y un despliegue reproducible con Docker Compose.

La principal conclusión es que la flexibilidad de un modelo de lenguaje y el
rigor factual de una base de datos pueden conciliarse cuando el papel del modelo
se limita a la comprensión y la redacción y todo dato se delega en la fuente
estructurada. El trabajo muestra, además, que en sistemas alimentados por
\emph{scraping} la calidad depende más de la normalización que de la extracción,
y que la metodología iterativa fue decisiva para incorporar decisiones de diseño
ausentes del plan inicial. \emph{autoAI} cumple así el objetivo planteado: dar a
un usuario no experto un acceso conversacional al catálogo de vehículos sin
renunciar a la veracidad de los datos.
